% TOP_PROBLEM: {"id": 30006894, "title": "Ivrii conjecture on periodic billiard orbits", "category_id": 11, "category": {"id": 11, "name": "dynamical_systems", "display_name": "Dynamical Systems", "description": "Problems about long-term behavior of deterministic systems, Hamiltonian dynamics, and periodic orbits.", "slug": "dynamical-systems", "order_index": 11, "created_at": "2026-07-31T15:26:25.671Z"}, "status": "open"}
% FIRST_POSTED: 2026-09-27
% !TeX program = lualatex
% ATTEMPT_STATUS: unresolved
% Model: GPT-6 Astra Ultra; continuation dated 27 September 2026.
\documentclass[11pt]{article}
\usepackage[margin=25mm]{geometry}
\usepackage{fontspec}
\setmainfont{Latin Modern Roman}
\usepackage{amsmath,amssymb,mathtools}
\usepackage{unicode-math}
\setmathfont{Latin Modern Math}
\usepackage{xurl,hyperref}
\hypersetup{colorlinks=true,urlcolor=blue}
\setcounter{secnumdepth}{0}
\setlength{\parindent}{0pt}
\setlength{\parskip}{5pt}
\setlength{\emergencystretch}{3em}
\begin{document}
\section*{\texorpdfstring{Ivrii conjecture on periodic billiard orbits}{Ivrii conjecture on periodic billiard orbits}}\label{30006894}
\subsection{Definitions and mathematical statement}
For a bounded connected smooth domain $\Omega\subset{\mathbb{R}}^n$, $n\ge2$, let $M=B^*\partial\Omega=\{(x,\eta):|\eta|<1\}$ be its open boundary cotangent unit ball bundle. Give it canonical Liouville measure $\mu=|\omega^{n-1}|/(n-1)!$. The billiard map $\beta$ follows the inward unit covector to its next boundary collision, reflects specularly, and records the tangential covector. Let $M_{\rm reg}$ consist of states with all positive and negative iterates well-defined and nongrazing. Ivrii's assertion is
\[
 \mu\!\left(\bigcup_{k\ge2}\{z\in M_{\rm reg}:\beta^kz=z\}\right)=0.
\]
Periods need not be least periods. Neither convexity nor simple connectivity is assumed, and singular trajectories are excluded rather than continued by an invented rule. The claim is zero phase-space measure, not finiteness or absence of periodic trajectories.
\par\smallskip\textit{Formulation and concept references:} \href{https://arxiv.org/html/1405.5990v4}{[S1]}.

\subsection{Short English statement}
In every bounded smooth Euclidean billiard domain, do regular periodic trajectories occupy zero Liouville measure in the boundary phase space?
\subsection{Research attempt}
\paragraph{The measure and the geometric category.}
The primary source [S1] studies a particular analytic planar
classification and its consequences for four reflections. Its
complexification and fixed-period results are not a proof for all
periods, all smooth boundaries or all dimensions. The catalogue
selects ordinary Euclidean specular billiards and the full regular
phase space. An example in a projective, outer or complex billiard
category would therefore not settle this formulation.

We first prove the two-collision assertion for every domain in the
question and the complete assertion for balls. We then give a precise
analytic reduction and show why smoothness and symplecticity alone
do not close its missing step. In local boundary coordinates
$(s,\eta)$ the Liouville measure is the ordinary symplectic volume
$ds\,d\eta$; smooth changes of boundary coordinates preserve its
measure class. Thus positive-dimensional families of periodic
trajectories can still have zero measure in the full
$2(n-1)$-dimensional phase space.

\paragraph{A finite-period reduction with the singular set kept out.}
For $k\geq2$, let $U_k$ be the open set of states for which the
first $k$ forward collisions are well-defined, have distinct
consecutive endpoints and are nongrazing. Smoothness of the boundary
and the implicit function theorem show that $\beta^k$ is smooth on
this open set, locally in a fixed itinerary of boundary charts.
Indeed the collision time is the first positive root of a boundary
defining function along a straight line; at a transverse hit its
time derivative is nonzero. Specular reflection is a smooth formula
in the incoming direction and normal. Repeating this argument gives
the claimed local regularity.

The periodic states in the question are contained in
\[
 \bigcup_{k\geq2}\{z\in U_k:\beta^kz=z\}.
 \tag{379.1}
\]
Thus it suffices to prove nullity for each fixed $k$. This is a
countable union, so no estimate uniform in $k$ is required merely
for the measure-zero conclusion. Conversely, checking finitely many
periods cannot establish (379.1). We do not assign a continuation
through tangency or a singular reflection in order to define an
iterate outside these regular sets.

\paragraph{All two-periodic states lie in a null section.}
Consider a regular orbit with two collisions at $x$ and $y$. The
outgoing velocity at $x$ is $(y-x)/|y-x|$, while the incoming
velocity there, after the return segment, is its negative. Let
$n_x$ be the outward normal. Specular reflection preserves the
tangential component of velocity. A vector and its negative have
the same tangential component only if that component is zero.
Thus the tangential covector at $x$ is $\eta=0$; the same is true
at $y$.

The zero section $\{(x,0):x\in\partial\Omega\}$ has dimension
$n-1$ inside a $2(n-1)$-dimensional manifold. Since $n\geq2$, its
Liouville measure is zero: in each local chart it is a product of
a boundary coordinate domain with a singleton in positive-dimensional
covector space, and Fubini applies. A countable chart cover proves
\[
 \mu\{z\in M_{\rm reg}:\beta^2z=z\}=0.
 \tag{379.2}
\]
The proof neither assumes convexity nor tries to count the normal
chords. There can be infinitely many such chords or a continuous
family, as in a ball, without violating (379.2).

\paragraph{A complete model: the ball in every dimension.}
By translation and scaling it suffices to take the unit ball. Write
$x$ for a boundary point and $v$ for the inward unit velocity. The
bivector $x\wedge v$ is conserved during a straight flight, since
$(x+tv)\wedge v=x\wedge v$. It is also unchanged at reflection:
the new velocity differs from the old one by a scalar multiple of
the radial normal $x$.

If $x\wedge v=0$, the orbit follows a diameter and has period two.
Otherwise the two-dimensional plane spanned by $x,v$ is invariant
under both the flights and reflections. To verify this directly,
the next collision point lies in that plane and the reflected
velocity is a linear combination of the incoming velocity and that
collision point. The entire orbit is therefore a circular billiard
inside this fixed plane.

Let $\rho=|\eta|\in[0,1)$ be the magnitude of the boundary
tangential velocity. It is conserved because it equals
$|x\wedge v|$ on the unit sphere. In the invariant circle, consecutive
collision points have central angle
\[
 \theta=2\arccos\rho\in(0,\pi],
 \tag{379.3}
\]
with a fixed orientation when $\rho>0$. For example $\rho=0$
gives the diameter angle $\pi$, while grazing would correspond
to $\rho=1$ and is excluded. The identity follows by resolving
the unit chord direction into its normal and tangential components
at one endpoint.

The state is periodic exactly when $\theta/(2\pi)$ is rational.
If $k\theta=2\pi m$, the collision point and its incoming and
outgoing directions repeat; conversely repetition of the state forces
this equality in its invariant circle. Thus all periodic states lie
in the union of fiber-radius levels
\[
 \bigcup_{\substack{k\geq2,\ 1\leq m\leq k/2}}
        \{(x,\eta):|\eta|=\cos(\pi m/k)\}.
 \tag{379.4}
\]
Every level has zero Lebesgue measure in the $(n-1)$-dimensional
covector fiber. For a positive radius it is a sphere of one lower
dimension; for radius zero it is a singleton. Fubini and the countable
union show that the full periodic set has Liouville measure zero.
This proves the selected assertion for every ball, in every $n\geq2$.

There are nevertheless many periodic trajectories in a ball. In
dimension two, rational angles are dense among all allowed angles,
so periodic states are dense. Density, infinitude and positive
measure are different properties. This model explicitly prevents
an argument that mistakes a dense periodic family for a counterexample.

\paragraph{An analytic zero-set lemma and its precise use.}
Let $F$ be a nonzero real-analytic scalar function on a connected
open subset of $\mathbb R^d$. Then its zero set has Lebesgue
measure zero. One proof keeps track of derivatives. At any zero
point, not every derivative can vanish, since otherwise the convergent
Taylor series and analytic continuation would make $F$ identically
zero. Choose a nonzero derivative of least total order $r\geq1$.
A derivative of order $r-1$ then vanishes at the point and has
nonzero gradient there. Its zero set is locally a smooth hypersurface
by the implicit function theorem. Stratifying according to the first
nonzero derivative covers the original zero set by countably many
such local hypersurfaces. Each is null, proving the assertion.

For a real-analytic billiard boundary, the collision-time argument
above is analytic on every regular itinerary chart. Hence
$\beta^k-\operatorname{id}$ is an analytic vector-valued map in
common coordinates near a possible fixed point. If one component
is not identically zero on a connected chart, its zero set, and
therefore the fixed-point set, is null there. This yields a conditional
reduction:
\[
 \begin{gathered}
 \text{If no regular analytic itinerary component has }
 \beta^k=\operatorname{id}\text{ on an open set,}\\
 \text{then the regular }k\text{-periodic set is null.}
 \end{gathered}
 \tag{379.5}
\]
One may use a countable coordinate cover, since the phase space is
second countable. Components without fixed points need no analysis;
near a fixed point the source and target of the iterate can be put
in the same coordinate chart.

The hypothesis excluding an identity germ is substantial. It is
not a consequence of the abstract fact that the billiard map is
analytic. Proving it for every period and every allowed table is
precisely a geometric issue left by this reduction. The source's
discussion of reflective analytic billiards makes this obstruction
visible rather than eliminating it for arbitrary periods.

\paragraph{Why a smooth zero-set argument is insufficient.}
A smooth function can vanish on a positive-measure set with empty
interior. To see this concretely, let $C$ be a closed nowhere-dense
subset of $[0,1]$ of positive measure, constructed by deleting a
countable collection of intervals of total length less than one.
Write the complementary intervals as $I_j=(a_j,b_j)$. On each interval
choose a nonnegative smooth bump $\psi_j$, positive inside and flat
at its endpoints. Multiply it by a coefficient so small that the
supremum of its first $j$ derivatives is at most $2^{-j}$.
The sum of these disjointly supported bumps, extended by zero on
$C$, is smooth; alternatively the same bounds give convergence of
every derivative of the series. It is positive off $C$ and zero
on $C$. Thus the analytic implication in (379.5) cannot be repeated
using only $C^\infty$ regularity and empty interior.

The chosen coefficients can also be required to control derivatives
near the finitely many early intervals, whose endpoint flatness
already gives continuity there. Hence the construction does not
depend on an unproved claim that an arbitrary sum of bumps is smooth.
It is an obstruction to a proposed method, not a constructed
Euclidean billiard with positive-measure periodic states.

\paragraph{Symplecticity does not exclude a positive fixed set.}
On a planar symplectic coordinate chart, take a smooth radial
Hamiltonian $H(q,p)=\phi(q^2+p^2)$ with $\phi$ constant for radii
near zero and outside a larger disc, and nonconstant in between.
Its Hamiltonian flow is the identity on the central disc and outside
the larger disc, and rotates intermediate circles by a radius-dependent
angle. Its time-one map preserves symplectic area and has an open
set of fixed points while being nonidentity elsewhere.

Thus measure preservation, symplecticity and smooth invertibility
do not by themselves imply the required nullity for a fixed-period
set. The billiard map belongs to a much more constrained class of
maps. Any proof based on fixed-point geometry must use that additional
Euclidean reflection structure, rather than applying a false general
principle to arbitrary symplectic maps.

There is a variational expression for the reflection law: for a
closed polygon of collision points $(x_1,\ldots,x_k)$, define
\[
 \mathcal L=\sum_{i=1}^k|x_{i+1}-x_i|,\qquad x_{k+1}=x_1.
\]
Its tangential derivative at $x_i$ is the tangential projection of
\[
 \frac{x_i-x_{i-1}}{|x_i-x_{i-1}|}
       +\frac{x_i-x_{i+1}}{|x_i-x_{i+1}|}.
\]
Vanishing is the reflection law for a nondegenerate itinerary with
the correct incoming and outgoing sides. However, a smooth function
can have degenerate critical manifolds. Counting equations and
variables in this length functional does not prove that the critical
set projects to a null set of initial states. The normal-chord and
ball calculations succeed because they identify extra equations
explicitly, not because of a generic dimension count.

\paragraph{Diagnostics and the remaining obstruction.}
The diagnostic iterates the exact circular collision-angle rule
for selected rational angles and checks the first return, including
the diameter case. It checks the angle formula by direct chord
geometry as well. These are consistency tests for the ball proof,
not evidence that finite-period enumeration exhausts a general table.

The proved results are nullity of the two-periodic set in every
selected domain, the full ball subcase, and the conditional analytic
reduction. In the analytic category the first gap is excluding
identity germs of every iterate for Euclidean billiards. In the
general smooth category even that exclusion alone does not replace
a measure-zero argument. The full assertion in the printed scope
remains unresolved by this attempt.
\subsection{Sources}
\begin{itemize}
\item[S1]Alexey Glutsyuk, On 4-reflective complex analytic planar billiards, arXiv:1405.5990v4 (17 December 2015); standard real billiard-map conventions as in Hassell-Zelditch, arXiv:math/0211140, Section 3.
\url{https://arxiv.org/html/1405.5990v4}.
\textit{Formulation source.}
\end{itemize}
\end{document}
